\section{Related work and the nearest-prior boundary}
\label{sec:related}

\paragraph{Intervention value at checkpoints.}
\Citet{calibrationnotcontrol2026}, \S3--4, executes continue, the benchmark
intervention and quit from each selected prefix, names \emph{intervention
advantage} as the decision quantity, and selects conservatively with a random
forest and a per-tree uncertainty penalty.  On interactive benchmarks it replays
the original trajectory, checks step agreement and \emph{discards} mismatches
(their \S3; Appendix B.1), so its reported match rate is a property of retained
rows, not an unconditional reconstruction guarantee.  That paper already
establishes prefix branching, action conditioning and conservative
action-value selection; neither branch labels, risk insufficiency nor an LCB
alone can carry a new contribution here.  We ask the timing question it does not
pose --- whether acting on one-step advantage spends the single verification too
early --- and compare fitted timing controllers under matched allowed budgets on
one shared tree.

\paragraph{Verification timing and budget deferral.}
\Citet{adeli2026bavar} studies strategic verification for long-running agents and
discusses verification timing, finite budgets, deferral for information and
reserving resources for later checks.  Our BAVAR-inspired scheduler is a
context-dependent reserve scheduler in that spirit, restricted to the signals a
single-verifier two-checkpoint task provides (\cref{sec:policy} lists the
signals it uses and its deviations).  A constant-reserve ablation is reported
separately; only the context-dependent scheduler is a primary comparison.
% UNVERIFIED CITATION (review M-2): the preprint page blocks automated access
% (checked 2026-09-27 and 2026-09-29), so the bibliographic data and the
% description above have not been checked at source. No section of the preprint
% is attributed until an author reads it.

\paragraph{Cost-aware information gathering.}
\Citet{calibratethenact2026} formalises cost--uncertainty tradeoffs and
cost-aware exploration in LLM agents (retrieval-augmented QA and a file-reading
coding task).  Our narrower contribution is the agent setting and the paired
counterfactual measurement design for original-root, action-dependent
two-checkpoint timing under an explicit replay-uncertainty contract, with
matched timing baselines.  It is a candidate contribution requiring
demonstration, not a new timing method.

\paragraph{Continuation-value and fitted-Q methodology.}
The proposed rule applies least-squares continuation-value estimation
\citep{longstaffschwartz2001} and fitted Q-iteration \citep{ernst2005} to agent
verification timing; the methodology is standard and is not claimed as new.

\paragraph{Lifecycle-wide interventions.}
\Citet{harnessr12026} reruns a full task batch after lifecycle hooks are
installed, keeps original successes in the batch, and reports that many valid
patches reduce batch success (\S3.1--3.2).  That is a transductive, whole-batch
rerun objective with lifecycle hooks, not a one-verification timing gate; we ask
when a proposed verification should be spent at a particular checkpoint under a
budget.

\paragraph{Replay and measurement preconditions.}
\Citet{causalagentreplay2026} re-execute trajectories under step interventions to
attribute failure, reporting an action-match rate, with side-effecting tools out
of scope.  \Citet{replaygap2026} fork trajectories by replaying logged actions
into a fresh container with a matched control arm and report reconstruction
rates.  Replay fidelity is measured there by action-match and reconstruction
rates.  The measurement contract this paper assumes --- per-layer restoration
and an explicit identified set for the intended effect under an unaudited
discrepancy --- is developed separately by the anonymous companion and applied
here; it is not re-derived.

\paragraph{Population information from repeated observations.}
\citet{tian2017populations} study population recovery from few repeated
binomial observations per entity. For B, this motivates separating independent
family count from within-root replay depth and considering information sharing
across training tasks. Such an estimator would be a proposed comparator, not a
result reported here. Our primary utility is bounded continuous accuracy, so a
binomial likelihood cannot simply be imposed on it. Nor does population
recovery alone identify context-specific continuation values or establish a
verification-timing advantage.

\paragraph{Positioning and contribution.}
The anonymous companion \citep{companion2026} establishes the population-scoped
$2p$ allowance used in \cref{prop:n1-population}; B applies it and does not
claim it as a new safety theorem.  The two papers are one program: the companion
fixes what a branch contrast identifies, and this paper asks when to spend a
verification budget under that constraint.  B's contribution is the agent
setting, the paired counterfactual measurement design, and convincing evidence
on it with explicitly scoped replay sensitivity.  It is \emph{not} a new timing
method, not an optimality claim, and not a claim that branch trees are the only
way to learn or evaluate a timing policy.

\paragraph{Ordinal sufficiency and selective prediction (displaced).}
The condition under which an outcome ranking coincides with an effect ranking is
established in budgeted triage \citep{okati2021}, in the uplift literature under
monotone treatment amplification \citep{fernandezloria2022}, and for allocation
at every budget fraction \citep{kamran2024}.  The exact ordinal result is prior
art; the implemented ridge advantage and action-value estimators are in fact the
same estimator, which is why that comparison is retained only as an equivalence
check (\cref{sec:b-equivalence}).  The fuller inherited ordinal, abstraction and
price-theoretic discussion is preserved in
\texttt{papers/paper\_b/archive/B1\_20260927/}.

\paragraph{What this paper does not claim.}
It does not claim that prediction and intervention value come apart; it does not
claim the ordinal-sufficiency theorem; it does not claim a new general
identification theorem for replay; it does not claim a new safety theorem from a
union bound; and it reports no real-agent result before the qualification and
pilot gates pass.
